% =====================================================================
% ABSTRACT + KEYWORDS
%
% El enunciado (6.2.1) pide TRES cosas y exige una en particular:
%   "El resumen se marca de forma explicita como PRELIMINAR, ya que
%    describe el problema, la motivacion y el enfoque propuesto sin
%    incorporar todavia resultados de ningun tipo."
% Por eso abre con "Preliminary abstract". No le quiten esa palabra
% hasta la Etapa 3, cuando ya haya resultados.
%
% CAMBIOS RESPECTO A LA VERSION ANTERIOR:
%  1. Dice PRELIMINARY, que es requisito de rubrica.
%  2. Antes decia "evaluating classical machine learning models and
%     tree-based ensembles"; ahora dice learning to rank, que es lo que
%     realmente propone la seccion de metodologia.
%  3. Antes prometia calibracion sin decir de donde salen las
%     probabilidades. LambdaMART devuelve scores sin escala, no
%     probabilidades. Ahora el abstract menciona el paso de calibracion
%     explicitamente.
%  4. "Palabras clave" -> "Index Terms" (el paper esta en ingles).
%
% PENDIENTE DEL EQUIPO: la metodologia dice que la metrica primaria es
% NDCG; este abstract nombra NDCG junto a top-1 y MRR. Si deciden que
% manda otra, ajusten aqui tambien.
% =====================================================================

\begin{abstract}
\textit{Preliminary abstract. This paper states the problem, the
motivation and the proposed approach; it reports no experimental results,
which are deferred to a later stage of the project.}

Horse race prediction has traditionally been framed either as an
independent binary classification problem over individual runners or as a
market-efficiency question in which tote odds serve as the main predictor.
Neither framing exploits the structure of the task: a race is a closed
field in which a known set of competitors is placed in a complete order.
We propose a learning-to-rank formulation over the full finishing order,
trained exclusively on information published before the starting gate
opens, using the Hong Kong Jockey Club records from June 1997 to August
2005. Admissibility is decided by an explicit column-by-column leakage
audit that removes post-race attributes such as finishing times, sectional
positions and payouts, and the evaluation uses a chronological split in
which every race remains an indivisible unit. Historical form is built
from strictly earlier races. Ranking scores are converted to win
probabilities by a calibration step fitted on the validation split, so
that discrimination and calibration can both be reported. Model quality is
assessed with ranking metrics together with race-level top-1 accuracy,
mean reciprocal rank, log loss and Brier score, and the whole ladder is
compared against the historical closing odds as an external market
benchmark.
\end{abstract}

\begin{IEEEkeywords}
machine learning, sports prediction, learning to rank, probability
calibration, data leakage
\end{IEEEkeywords}